\documentclass[a4paper,12pt]{article}

\usepackage[T2A]{fontenc}
\usepackage[utf8]{inputenc}
\usepackage[russian]{babel}

\usepackage{amsmath,amssymb,amsthm,mathtools}
\usepackage{helvet}
\renewcommand{\familydefault}{\sfdefault}
\usepackage{icomma}
\usepackage[a4paper,margin=20mm,headheight=15pt]{geometry}
\usepackage{microtype}

\usepackage{tikz}
\usetikzlibrary{calc}
\usetikzlibrary{arrows.meta,positioning,shapes.geometric}

\usepackage{tkz-euclide}
\usepackage{pgfplots}
\pgfplotsset{compat=1.18}

\usepackage{tabularx,booktabs,longtable}
\usepackage{enumitem}
\usepackage{hyperref}
\usepackage{parskip}
\usepackage{fancyhdr}
\usepackage{setspace}
\usepackage{xcolor}

% ------------------------------------------------
% Цветовая палитра
% ------------------------------------------------
\definecolor{accent}{HTML}{008080}
\definecolor{darkaccent}{HTML}{004D4D}
\definecolor{textgray}{HTML}{333333}
\definecolor{lightgray}{HTML}{E7ECEC}

\color{textgray}
\onehalfspacing

\hypersetup{
    colorlinks=true,
    linkcolor=darkaccent,
    urlcolor=darkaccent
}

\setlength{\parindent}{0pt}
\setlist[itemize]{leftmargin=6mm,itemsep=2pt,topsep=3pt}
\setlist[enumerate]{leftmargin=7mm,itemsep=4pt,topsep=4pt}

% ------------------------------------------------
% Колонтитулы
% ------------------------------------------------
\pagestyle{fancy}
\fancyhf{}
\fancyhead[L]{\small Надя Клименко}
\fancyhead[R]{\small Математика, 6 класс}
\fancyfoot[C]{\small\thepage}
\renewcommand{\headrulewidth}{0.3pt}
\renewcommand{\footrulewidth}{0pt}

% ------------------------------------------------
% Заголовки
% ------------------------------------------------
\newcommand{\sectionline}{%
    \par\vspace{2mm}
    {\color{accent}\rule{\linewidth}{0.7pt}}
    \vspace{2mm}
}

\newcommand{\checkitem}[1]{%
    \item[\textcolor{accent}{$\square$}] #1
}

% ------------------------------------------------
% TikZ-стили для блок-схемы
% ------------------------------------------------
\tikzset{
    flowstart/.style={
        ellipse,
        draw=accent,
        line width=0.9pt,
        minimum width=34mm,
        minimum height=11mm,
        align=center,
        font=\small
    },
    flowaction/.style={
        rectangle,
        rounded corners=2mm,
        draw=accent,
        line width=0.9pt,
        text width=48mm,
        minimum height=13mm,
        align=center,
        inner sep=4mm,
        font=\small
    },
    flowdecision/.style={
        diamond,
        aspect=2.15,
        draw=accent,
        line width=0.9pt,
        text width=37mm,
        align=center,
        inner sep=2mm,
        font=\small
    },
    flowarrow/.style={
        -{Latex[length=3mm,width=2mm]},
        line width=0.9pt,
        draw=darkaccent
    },
    flowlabel/.style={
        font=\small,
        fill=white,
        inner sep=1pt
    }
}

\begin{document}

% =========================================================
% ТИТУЛЬНЫЙ ЛИСТ
% =========================================================
\begin{titlepage}
\thispagestyle{empty}

\begin{center}

\vspace*{22mm}

{\color{darkaccent}\Large\bfseries Чек-лист}

\vspace{8mm}

{\Huge\bfseries
Уравнения, проценты, НОД и НОК
}

\vspace{5mm}

{\Large
Ключевые методы и типовые задачи
}

\vspace{20mm}

{\large
Надя Клименко
}

\vspace{3mm}

{\large
6 класс
}

\vfill

{\large
\today
}

\vspace{5mm}

{\large
Лёвин Артём Александрович
}

\vspace{2mm}

{\large
эксклюзивно для Нади Клименко
}

\vspace{18mm}

\begin{tikzpicture}[remember picture,overlay]
    \draw[
        color=accent!65,
        line width=1.4pt
    ]
    ([xshift=18mm,yshift=17mm]current page.south west)
    .. controls
    ([xshift=65mm,yshift=32mm]current page.south west)
    and
    ([xshift=-70mm,yshift=8mm]current page.south east)
    ..
    ([xshift=-18mm,yshift=20mm]current page.south east);
\end{tikzpicture}

\end{center}
\end{titlepage}

% =========================================================
% ОСНОВНОЙ ЧЕК-ЛИСТ
% =========================================================
\section*{\color{darkaccent}1. Уравнения: восстанавливаем неизвестный компонент}
\sectionline

\begin{itemize}
    \checkitem{Понимаю смысл равенства
    \[
        \frac{A}{B}=C
        \quad\Longrightarrow\quad
        B=\frac{A}{C},
        \qquad C\ne0.
    \]}

    \checkitem{Перед решением дробного уравнения проверяю, при каких значениях знаменатель не обращается в ноль.}

    \checkitem{Если неизвестное находится в знаменателе, сначала выражаю весь знаменатель, затем нахожу неизвестное.}

    \checkitem{Если неизвестное является уменьшаемым или вычитаемым, использую смысл действия вычитания, а не механический ``перенос''.}
\end{itemize}

\subsection*{Пример 1. Неизвестное в знаменателе}

Решим
\[
    \frac{6}{x-1}=2.
\]

Сначала выпишем ограничение:
\[
    x-1\ne0,
    \qquad
    x\ne1.
\]

Из равенства
\[
    \frac{6}{x-1}=2
\]
получаем
\[
    x-1=\frac{6}{2}=3.
\]

Следовательно,
\[
    x=4.
\]

Проверка:
\[
    \frac{6}{4-1}=\frac63=2.
\]

Ответ:
\[
    \boxed{x=4}.
\]

\subsection*{Пример 2. Неизвестное в числителе}

\[
    \frac{x-1}{2}=3.
\]

Умножаем обе части на \(2\):
\[
    x-1=6,
\]
откуда
\[
    x=7.
\]

\subsection*{Пример 3. Неизвестное вычитаемое}

\[
    10-x=6.
\]

Полезно вспомнить простое равенство:
\[
    10-2=8.
\]

Чтобы найти вычитаемое, нужно из уменьшаемого вычесть разность:
\[
    x=10-6=4.
\]

\textbf{Типичная ошибка:}
из \(10-x=6\) нельзя получать \(x=10+6\).

\textbf{Контроль смысла:}
\[
    10-4=6.
\]

% =========================================================
\section*{\color{darkaccent}2. Проценты: три основных типа задач}
\sectionline

\subsection*{2.1. Найти \(p\%\) от числа}

Чтобы найти \(p\%\) от числа \(A\), используйте
\[
    A\cdot\frac{p}{100}.
\]

Например,
\[
    20\%\text{ от }500
    =
    500\cdot\frac{20}{100}
    =
    100.
\]

\subsection*{2.2. Найти число по известному проценту}

Если \(p\%\) некоторого числа равны \(B\), то всё число равно
\[
    B\cdot\frac{100}{p}.
\]

Например, если \(20\%\) числа равны \(500\), то
\[
    500\cdot\frac{100}{20}
    =
    2500.
\]

\subsection*{2.3. Найти, сколько процентов одно число составляет от другого}

Если нужно определить, сколько процентов число \(A\) составляет от числа \(B\), используйте
\[
    \frac{A}{B}\cdot100\%.
\]

Например,
\[
    \frac{20}{100}\cdot100\%=20\%.
\]

\subsection*{2.4. Скидка}

Если цена товара равна \(A\), а скидка составляет \(p\%\), можно действовать двумя способами.

\textbf{Способ 1.}
Сначала найти величину скидки:
\[
    A\cdot\frac{p}{100},
\]
затем вычесть её из исходной цены.

\textbf{Способ 2.}
Сразу найти оставшийся процент:
\[
    100\%-p\%.
\]

Например, при скидке \(20\%\) остаётся заплатить \(80\%\) цены.

Для футболки стоимостью \(3000\) рублей:
\[
    3000\cdot\frac{80}{100}=2400.
\]

Ответ:
\[
    \boxed{2400\text{ руб.}}
\]

\textbf{Типичная ошибка:}
скидка \(20\%\) означает, что покупатель платит \(80\%\) исходной цены.

% =========================================================
\section*{\color{darkaccent}3. НОД: когда объекты нужно разделить на одинаковые группы}
\sectionline

\textbf{Наибольший общий делитель} чисел \(a\) и \(b\) --- это наибольшее натуральное число, на которое делятся и \(a\), и \(b\).

Обозначение:
\[
    \operatorname{НОД}(a,b).
\]

Если чисел несколько:
\[
    \operatorname{НОД}(a,b,c).
\]

\begin{itemize}
    \checkitem{Распознаю НОД по словам: ``разделить'', ``поровну'', ``одинаковые группы'', ``одинаковые комплекты'', ``наибольшее возможное количество групп''.}

    \checkitem{Умею разложить числа на простые множители.}

    \checkitem{Для НОД беру только те простые множители, которые встречаются в разложении каждого числа, причём в наименьшей необходимой степени.}
\end{itemize}

\subsection*{Пример}

После предварительной раздачи осталось:
\[
    252\text{ тетради},\qquad
    392\text{ ручки},\qquad
    504\text{ карандаша}.
\]

Нужно составить максимально возможное число одинаковых комплектов.

Разложим числа:
\[
    252=2^2\cdot3^2\cdot7,
\]
\[
    392=2^3\cdot7^2,
\]
\[
    504=2^3\cdot3^2\cdot7.
\]

Общие множители:
\[
    2^2\cdot7.
\]

Следовательно,
\[
    \operatorname{НОД}(252,392,504)
    =
    2^2\cdot7
    =
    28.
\]

Можно составить
\[
    \boxed{28}
\]
одинаковых комплектов.

В каждом комплекте:
\[
    \frac{252}{28}=9
    \quad\text{тетрадей},
\]
\[
    \frac{392}{28}=14
    \quad\text{ручек},
\]
\[
    \frac{504}{28}=18
    \quad\text{карандашей}.
\]

\textbf{Важно.}
НОД не обязан быть строго меньше каждого исходного числа.

Например,
\[
    \operatorname{НОД}(6,12)=6.
\]

% =========================================================
\section*{\color{darkaccent}4. НОК: когда нужно найти первое общее совпадение}
\sectionline

\textbf{Наименьшее общее кратное} чисел \(a\) и \(b\) --- наименьшее натуральное число, которое делится и на \(a\), и на \(b\).

Обозначение:
\[
    \operatorname{НОК}(a,b).
\]

\begin{itemize}
    \checkitem{Распознаю НОК по словам: ``одновременно'', ``снова совпадут'', ``первая общая точка'', ``наименьшая общая сумма'', ``через сколько повторится''.}

    \checkitem{Для НОК беру множители одного числа и добавляю недостающие множители остальных чисел.}

    \checkitem{Проверяю, что найденное число делится на каждое исходное число без остатка.}
\end{itemize}

\subsection*{Пример 1. Одинаковая наименьшая сумма}

Один пирожок стоит \(45\) рублей, другой --- \(60\) рублей.

Нужно найти наименьшую сумму, которую можно потратить в обоих случаях, купив целое число пирожков.

Разложения:
\[
    45=3^2\cdot5,
\]
\[
    60=2^2\cdot3\cdot5.
\]

Для НОК необходимы множители
\[
    2^2,\qquad 3^2,\qquad 5.
\]

Поэтому
\[
    \operatorname{НОК}(45,60)
    =
    2^2\cdot3^2\cdot5
    =
    180.
\]

На \(180\) рублей можно купить:
\[
    \frac{180}{45}=4
\]
пирожка первого вида и
\[
    \frac{180}{60}=3
\]
пирожка второго вида.

\subsection*{Пример 2. НОК чисел \(15\) и \(24\)}

\[
    15=3\cdot5,
\qquad
    24=2^3\cdot3.
\]

Следовательно,
\[
    \operatorname{НОК}(15,24)
    =
    2^3\cdot3\cdot5
    =
    120.
\]

Число \(240\) тоже делится и на \(15\), и на \(24\), однако
\[
    240>120,
\]
поэтому оно не является \emph{наименьшим} общим кратным.

\textbf{Важно.}
НОК не обязан быть строго больше каждого исходного числа.

Например,
\[
    \operatorname{НОК}(6,12)=12.
\]

% =========================================================
\section*{\color{darkaccent}5. Как различать НОД и НОК}
\sectionline

\begin{center}
\renewcommand{\arraystretch}{1.35}
\begin{tabularx}{\linewidth}{@{}>{\bfseries}p{0.19\linewidth}XX@{}}
\toprule
 & НОД & НОК \\
\midrule
Главный смысл
&
Число, которое \emph{делит} данные числа
&
Число, которое \emph{делится} на данные числа
\\

Типичный сюжет
&
Разделить предметы на одинаковые группы
&
Найти первое общее совпадение или общую наименьшую величину
\\

При разложении
&
Берём только общие множители
&
Берём все необходимые множители
\\

Пример
&
\(\operatorname{НОД}(15,24)=3\)
&
\(\operatorname{НОК}(15,24)=120\)
\\
\bottomrule
\end{tabularx}
\end{center}

\textbf{Короткая опора для памяти:}
\[
    \boxed{\text{НОД делит числа}}
    \qquad
    \boxed{\text{НОК делится на числа}}.
\]


% =========================================================
% ТРЕНИРОВОЧНЫЙ БЛОК
% =========================================================
\section*{\color{darkaccent}Тренировочный блок. Путь А}
\sectionline

\textbf{Задание 1. Уравнения}

Решите уравнения.

\begin{enumerate}[label=\alph*)]
    \item
    \[
        \frac{18}{x-2}=6;
    \]

    \item
    \[
        \frac{x+3}{5}=4;
    \]

    \item
    \[
        17-x=9.
    \]
\end{enumerate}

\vspace{7mm}

\textbf{Задание 2. Проценты}

Куртка стоила \(4800\) рублей.

Сначала цену снизили на \(15\%\).

Найдите стоимость куртки после скидки.

\vspace{7mm}

\textbf{Задание 3. НОД и НОК}

Для школьного мероприятия подготовили \(168\) тетрадей и \(252\) карандаша.

\begin{enumerate}[label=\alph*)]
    \item Какое наибольшее число одинаковых наборов можно составить так, чтобы использовать все тетради и карандаши?

    \item Сколько тетрадей и сколько карандашей будет в одном наборе?

    \item Два автоматических устройства включаются через каждые \(18\) минут и \(24\) минуты соответственно. Сейчас они включились одновременно. Через сколько минут они впервые снова включатся одновременно?
\end{enumerate}

\clearpage

% =========================================================
% ОТВЕТЫ
% =========================================================
\section*{\color{darkaccent}Ответы}
\sectionline

\renewcommand{\arraystretch}{1.45}

\begin{center}
\begin{tabularx}{\linewidth}{@{}p{0.16\linewidth}X@{}}
\toprule
\textbf{Задание} & \textbf{Ответ} \\
\midrule

1а &
\[
x=5
\]
\\

1б &
\[
x=17
\]
\\

1в &
\[
x=8
\]
\\

2 &
\[
4080\text{ руб.}
\]
\\

3а &
\[
84\text{ набора}
\]
\\

3б &
\[
2\text{ тетради и }3\text{ карандаша}
\]
\\

3в &
\[
72\text{ минуты}
\]
\\

\bottomrule
\end{tabularx}
\end{center}

\vfill

{\small
\textbf{Самопроверка перед завершением работы:}
проверяйте ограничения в дробных уравнениях; различайте исходное число и его процент; в задачах на одинаковые группы ищите НОД; в задачах на первое общее совпадение ищите НОК.
}

\end{document}